\chapter{Answers to the questions}\label{app:answers}

\paragraph{Question~\ref{q:oddcycle}.} $\val(G)=1-1/L$. A labeling that satisfies every edge has $\lambda(v_{i+1})=\lambda(v_i)+1$ for all $i$, and summing around the cycle gives $\lambda(v_1)=\lambda(v_1)+L=\lambda(v_1)+1$, as in Figure~\ref{fig:ugexample}. So every labeling fails at least one of the $L$ equally weighted edges, and $\val(G)\le1-1/L$. The labeling $\lambda(v_i)=i\bmod2$ satisfies the edges $(v_i,v_{i+1})$ for $i<L$, so $\val(G)\ge1-1/L$.

\paragraph{Question~\ref{q:propagation}.} In a connected component, fix a spanning tree and a label at its root. A permutation constraint on a tree edge leaves exactly one label at the far endpoint that satisfies the edge. So each of the $|K|$ root labels determines one labeling of the component, and checking each such labeling on every edge takes polynomial time. A 2-to-2 constraint relates each label at one endpoint to two labels at the other. Each root label then leaves two choices at every tree edge, so the number of labelings that satisfy the tree edges is exponential in the number of tree edges, and enumerating these labelings takes exponential time.

\paragraph{Question~\ref{q:evvexact}.} By Lemma~\ref{lem:evv}, it suffices to find a labeling with value $1-\bigl(1-\OPT(\Theta)\bigr)/3$. Take an assignment $\mathbf x$ that attains $\OPT(\Theta)$, and label each variable by its value. Label each equation $\theta$ by the restriction of $\mathbf x$ to the variables of $\theta$ if $\mathbf x$ satisfies $\theta$, and otherwise by that restriction with one variable flipped, so that the label satisfies $\theta$. Given $\theta$, the random edge then fails with probability $0$ in the first case, and with probability exactly $1/3$ in the second, since the label differs from $\mathbf x$ at exactly one of the three variables. So this labeling has value $1-\tfrac13\Prob_{\theta\sim p}[\mathbf x\text{ violates }\theta]=1-\bigl(1-\OPT(\Theta)\bigr)/3$.

\paragraph{Question~\ref{q:pibar}.} With coordinates $(z_0,x_3)$ on $L_O$, restriction to $x_3$ is $z\mapsto(z_0,\,bz_0+\phi_1+\phi_2)$, with matrix $\begin{pmatrix}1&0&0\\ b&1&1\end{pmatrix}$. With coordinates $(z_0,x_1)$, restriction to $x_1$ is $z\mapsto(z_0,\phi_1)$, with matrix $\begin{pmatrix}1&0&0\\0&1&0\end{pmatrix}$. Both keep the homogeneous coordinate. On the bits alone, restriction to $x_1$ is the linear map $(\phi_1,\phi_2)\mapsto\phi_1$, while restriction to $x_3$ is $(\phi_1,\phi_2)\mapsto b+\phi_1+\phi_2$, which is linear only when $b=0$.

\paragraph{Question~\ref{q:homog}.} Every label $z$ has $z_0=1$, so every encoding satisfies $f_z(M^{+h})=Mz+hz_0=f_z(M)+h$. Folding imposes this identity, which all encodings share, on every labeling. Adding $h$ to the column at a coordinate $i\ne0$ gives $f_z(M+he_i^\top)=f_z(M)+hz_i$, and the change $hz_i$ depends on the label. So this shift changes different encodings by different amounts, and folding cannot impose it.

\paragraph{Question~\ref{q:disagree}.} The probability is $2^{-\ell}$. Since $(P'\bar\pi_i)z_i=P'(\bar\pi_iz_i)$, the two values are $P'(\bar\pi_1z_1)$ and $P'(\bar\pi_2z_2)$, and they agree exactly when $P'(\bar\pi_1z_1-\bar\pi_2z_2)=0$. Since $\bar\pi_1z_1-\bar\pi_2z_2\ne0$ and $P'$ is a uniform linear map $L_O\to K$, this vector $P'(\bar\pi_1z_1-\bar\pi_2z_2)$ is uniform on $K$. So encodings of labels with different restrictions disagree on a $1-2^{-\ell}$ fraction of the tables on $O$.

\paragraph{Question~\ref{q:rowslices}.} Fix $d\in\F_2^\ell\setminus\{0\}$, write $t'=M^\top d$, and let $f(M)=Mt'$. On the row slice $M^\top d=t'$, the function $f$ is the encoding $f_{t'}$.

\emph{$f$ passes the equality test with probability at least $1/4$.} If $a^\top d=0$ and $l^\top t'=0$, then $(M+al^\top)^\top d=M^\top d=t'$, so the step stays in its row slice, and $f(M+al^\top)=Mt'+a(l^\top t')=f(M)$. Given $M$, the two events are independent, with probabilities $1/2$ and at least $1/2$.

\emph{$f$ agrees with every target on a small fraction of every slice cut out by column equations alone.} Fix such a slice, $Mq_j=t_j$ for $j=1,\dots,r$, and a target $Mz+u$. On the slice, $t'$ is uniform on a coset of the space $\{y\in\F_2^m:y^\top q_j=0\text{ for all }j\}$, whose dimension is at least $m-r$. The event $f(M)=Mz+u$ reads $M(t'-z)=u$. The vector $t'-z$ lies in $\Span\{q_j\}$ with probability at most $2^r\cdot2^{-(m-r)}$. Otherwise, given $t'$, the vector $M(t'-z)$ is uniform on a coset of the hyperplane $\{y\in K:d^\top y=0\}$, so $M(t'-z)=u$ has probability at most $2^{1-\ell}$. Hence the agreement on the slice is at most $2^{1-\ell}+2^{2r-m}$.

The function of Example~\ref{ex:zoomout} has small agreement on every slice cut out by row equations alone. Write the row equations as $AM=S$, with $s=\rank A\le r$, and fix a target $Mz+u$ with $z=z'+\nu q$ as in Example~\ref{ex:zoomout}. Given the nonempty row fiber, $t=Mq$ is uniform on a coset of $\ker A$, which has $2^{\ell-s}$ elements. At most one $t$ has $z_t=z'$, and that $t$ contributes at most $2^{-(\ell-s)}$. For the other $t$, the vector $z_t-z'$ is independent of $q$, so given $Mq=t$ the vector $M(z_t-z')$ is uniform on a coset of $\ker A$ and equals $\nu t+u$ with probability at most $2^{-(\ell-s)}$. So the agreement is at most $2^{1+r-\ell}$. So a theorem of the form of Theorem~\ref{thm:inverse} must allow both row and column equations.

\paragraph{Question~\ref{q:beta}.} \emph{$\beta=1$.} Every entry is replaced. Given $U$ and the positions of the replacing variables, the pullbacks $M_O\bar\pi$ form a subspace of $\Hom(L_U,K)$ of dimension $\ell(1+k)$, because $\bar\pi$ is onto and $L_O$ has dimension $1+k$. There are $3^k$ choices of positions, so the law of $M_O\bar\pi$ is supported on at most $3^k2^{\ell(1+k)}$ of the $2^{\ell(1+2k)}$ tables. Its total variation distance from the uniform law is therefore at least $1-(3\cdot2^{-\ell})^k$. So the conclusion of Lemma~\ref{lem:tv} fails, and the proof of Lemma~\ref{lem:transfer}, which carries good advice from uniform tables to pullbacks, no longer applies. Steps~1, 2, and~5 still hold, and so does Lemma~\ref{lem:goodadvice}, which concerns uniform tables. Step~4 fails.

\emph{$\beta=1/k$.} Now $k\beta^2=1/k\to0$, so Lemma~\ref{lem:tv} holds, but Step~5 fails. With probability $(1-1/k)^k\ge1/4$ for $k\ge2$, no entry is replaced. Then $O=U$, the map $\bar\pi$ is the identity, both provers receive the same input, and two copies of one deterministic strategy succeed. So no bound that tends to $0$ as $k$ grows can hold. Correspondingly, the bound of Lemma~\ref{lem:clean} becomes $\bigl(1-2^{-r_{\mathrm s}}/(3600k)\bigr)^k$, which tends to $\exp(-2^{-r_{\mathrm s}}/3600)$.

\paragraph{Question~\ref{q:erasure}.} A constant would create agreement instead of erasing it. A function $f'$ that took a constant value $y$ on $\mathcal B$ would agree with the target $M\mapsto M0+y$ at every table of $\mathcal B$. Every slice inside the row fiber of a good triple would then carry agreement $1$, and Theorem~\ref{thm:inverse} would not be contradicted. Each new value equals a fixed target with probability exactly $2^{-\ell}$, so the expected agreement of the new values with a target on a slice is at most $2^{-\ell}<\alpha/8$. Hoeffding's inequality, with a union bound over the at most $2^{(2r_{\mathrm s}+1)(\ell+m)}$ descriptions of slices and targets, keeps the contribution of the resampled values below $\alpha/4$ on every slice, for every target.

\paragraph{Question~\ref{q:mu}.} For $a\in K$, the conditions $g_i(a)=0$ for all $i$ say that $a$ lies in the kernel of the linear map $K\to\F_2^m$, $a\mapsto(g_0(a),\dots,g_{m-1}(a))$. The rank of this map is $\dim\Span\{g_i|_K\}=\rho(\boldsymbol g)$, so its kernel has dimension $\ell-\rho(\boldsymbol g)$. A uniform $a\in K$ lies in it with probability $2^{\ell-\rho(\boldsymbol g)}/2^\ell=2^{-\rho(\boldsymbol g)}$.

\paragraph{Question~\ref{q:point8}.} Nothing useful. With $q=0.8$ the computation gives $w\ge8q-7=-0.6$, which is no bound. The trivial character contributes $2^{-\ell}$ to $w$, and $2^{-\ell}\to0$ as $\ell$ grows. So the comparison does not give the acceptance bound independent of $\ell$ that Theorem~\ref{thm:inverse} needs.

\paragraph{Question~\ref{q:holenstein}.} The proof needs $n$ with $\varepsilon/n<\Delta$ and $\zeta^{n/\ell(\varepsilon/n)}\le\delta$, that is,
\[
\frac{n}{\ell(\varepsilon/n)}\ \ge\ \frac{\log(1/\delta)}{\log(1/\zeta)}.
\]
With $\varepsilon'=\varepsilon/n$, the left side is $n/\ell(\varepsilon')=\varepsilon/\bigl(\varepsilon'\ell(\varepsilon')\bigr)$. So such $n$ exists for every $\delta$ exactly when $\varepsilon'\ell(\varepsilon')$ has infimum $0$ along $\varepsilon'=\varepsilon/n$. The infimum is $0$ when $\ell(\varepsilon')=O(\log(1/\varepsilon'))$ or $\ell(\varepsilon')=O(\varepsilon'^{-1/2})$, and positive when $\varepsilon'\ell(\varepsilon')$ is bounded below by a positive constant.

\paragraph{Question~\ref{q:subdivision}.} Let $\lambda'$ label the new game, and let $\lambda$ be the restriction of $\lambda'$ to the vertices of $H$. The four translations of a path compose to the translation of $e$, so if $\lambda$ fails $e$, then $\lambda'$ fails at least one of the four path edges. Hence $\val_{\lambda'}\le1-\tfrac14\Prob_{e\sim p}[\lambda\text{ fails }e]\le1-(1-\val(H))/4$. Conversely, take a labeling $\lambda$ of $H$ with $\val_\lambda(H)=\val(H)$, and on each path choose the labels of $x_1$, $x_2$, and $x_3$ so that the first three path edges are satisfied. The fourth edge is then satisfied exactly when $\lambda$ satisfies $e$, so the extended labeling has value $1-(1-\val(H))/4$.
